% answers_short.tex —— 2013 年·湖北卷·理科 · 简短答案速查
% 由 tools/make_short_answers.py 自 answers.tex 的 \gksolhead/\gkanswer 提取，
% 并与 2013/2013_数学真题_湖北理_答案要点.md 逐题比对后生成，勿手改（改源 MD 后重跑）。

\begin{gknotice}
\textbf{使用说明}：本节只给最终答案便于快速核对，完整推导见解析版。答案为编写者独立推导，\textbf{非官方标准答案}。
\end{gknotice}

\gkpart{选择题}
\begin{center}\begin{tabular}{ccccccccccc}
\toprule
题号 & 1 & 2 & 3 & 4 & 5 & 6 & 7 & 8 & 9 & 10 \\
\midrule
答案 & D & C & A & B & D & A & C & C & B & D \\
\bottomrule
\end{tabular}\end{center}

\gkpart{填空题}
\begin{itemize}[leftmargin=2.6em]
  \item[\textbf{11.}] (1) $0.0044$。由 $50(0.0024+0.0036+0.0060+x+0.0024+0.0012)=1$ 解得；(2) $70$。$50(0.0036+0.0060+0.0044)\times100=70$
  \item[\textbf{12.}] $5$。$a$ 依次取 $10,5,16,8,4$，$i$ 依次为 $1,2,3,4,5$，$a=4$ 时输出 $i=5$。
  \item[\textbf{13.}] $\dfrac{3\sqrt{14}}{7}$。由柯西不等式取等得 $(x,y,z)=\dfrac1{\sqrt{14}}(1,2,3)$。
  \item[\textbf{14.}] $1000$。$N(n,k)=\dfrac{k-2}{2}n^2-\dfrac{k-4}{2}n$，故 $N(10,24)=10(22\cdot10-20)/2$。
  \item[\textbf{15.}] $8$。设半径为 $R$，则 $OD=\dfrac R3$，$OE=\dfrac R9$，$CE=\dfrac{8R}{9}$。
  \item[\textbf{16.}] $\dfrac{\sqrt6}{3}$。直线 $x+y=m$ 与圆相切得 $m=\sqrt2b$，过焦点得 $c=m$，故 $e=\dfrac{c}{a}=\dfrac{\sqrt2b}{\sqrt3b}$。
\end{itemize}

\gkpart{解答题}
\begin{itemize}[leftmargin=2.6em]
  \item[\textbf{17.}] (1) $A=\dfrac\pi3$。由 $\cos2A+3\cos A=1$ 得 $(2\cos A-1)(\cos A+2)=0$；(2) $\dfrac57$。$bc=20$，$c=4$，$a^2=21$，$\sin B\sin C=\dfrac{bc\sin^2A}{a^2}=\dfrac57$
  \item[\textbf{18.}] (1) $a_n=5\cdot3^{n-2}$ 或 $a_n=5\cdot(-1)^n$。由 $a_2=5$ 及 $|1-q|=2$ 得 $q=3$ 或 $q=-1$；(2) 不存在。$q=3$ 时部分和 $<0.9$；$q=-1$ 时部分和仅为 $0$ 或 $-\dfrac15$，均小于 $1$
  \item[\textbf{19.}] (1) $l\parallel$ 平面 $PAC$（证明见解析）。由 $EF\parallel AC$ 得 $EF\parallel$ 平面 $ABC$，从而 $EF\parallel l$；(2) 证明见解析。以 $O$ 为原点建系，算出 $\sin\theta$、$\sin\alpha$、$\sin\beta$ 后相乘即得
  \item[\textbf{20.}] (1) $0.9772$。$P_0=P(X\leqslant\mu+2\sigma)=0.5+\dfrac{0.9544}{2}$；(2) $A$ 型车 $5$ 辆、$B$ 型车 $12$ 辆，最小营运成本 $36800$ 元。由 $36x+60y\geqslant900$ 得 $5y\geqslant75-3x$，从而 $z\geqslant160x+36000\geqslant36800$
  \item[\textbf{21.}] (1) $\lambda=1+\sqrt2$。由 $S_1=\dfrac12a(m+n)$、$S_2=\dfrac12a(m-n)$ 得 $m^2-2mn-n^2=0$；(2) 当 $\lambda>1+\sqrt2$ 时存在；当 $1<\lambda\leqslant1+\sqrt2$ 时不存在。由 $\dfrac{S_1}{S_2}=\dfrac{r_m+r_n}{r_m-r_n}$ 及 $\dfrac{r_m}{r_n}\in(1,\lambda)$ 判断
  \item[\textbf{22.}] (1) $0$。$f'(x)=(r+1)\left[(1+x)^r-1\right]$，在 $x=0$ 处取得最小值；(2) 证明见解析。在（1）中分别取 $x=\dfrac1n$ 与 $x=-\dfrac1n$；(3) $211$。由（2）夹逼得 $210.225<S<210.9$
\end{itemize}
